% =============================================================================
% network-protocol.tex  —  Specifica Protocollo di Rete AM54
% Compilare con pdfLaTeX
% =============================================================================
\documentclass[a4paper,12pt]{article}

\usepackage[utf8]{inputenc}
\usepackage[T1]{fontenc}
\usepackage[italian]{babel}
\usepackage{amsmath}

\usepackage{geometry}
\geometry{margin=2.5cm, headheight=14pt, includehead, includefoot}

\usepackage{microtype}
\usepackage{graphicx}
\usepackage{xcolor}
\usepackage{tikz}
\usetikzlibrary{arrows.meta, calc, positioning, decorations.pathmorphing}

\usepackage{booktabs}
\usepackage{tabularx}
\usepackage{array}
\usepackage{float}
\usepackage{caption}
\usepackage{parskip}
\usepackage{enumitem}
\usepackage{fancyhdr}
\usepackage{lastpage}
\usepackage{hyperref}

\hypersetup{colorlinks=true, linkcolor=blue, urlcolor=blue}

% --------------- Intestazioni ---------------
\pagestyle{fancy}
\fancyhf{}
\fancyhead[L]{\small Mesos --- Specifica Protocollo di Rete}
\fancyhead[R]{\small Pagina \thepage/\pageref{LastPage}}
\fancyfoot[C]{\small \texttt{IS26-AM54} --- Cacioppo, Carbone, Buca, Botton}
\renewcommand{\headrulewidth}{0.4pt}
\renewcommand{\footrulewidth}{0.4pt}
\fancypagestyle{plain}{\fancyhf{}\renewcommand{\headrulewidth}{0pt}\renewcommand{\footrulewidth}{0pt}}

% --------------- Colori ---------------
\definecolor{clGreen}{HTML}{E8F5E9}
\definecolor{srvOrange}{HTML}{FFF3E0}
\definecolor{mdlRed}{HTML}{FFEBEE}
\definecolor{sysGray}{HTML}{F5F5F5}
\definecolor{noteBlue}{HTML}{E3F2FD}

% --------------- TikZ stili globali ---------------
\tikzset{
  hdr/.style={draw, rectangle, rounded corners=3pt,
              minimum width=2.1cm, minimum height=0.55cm,
              font=\scriptsize\ttfamily, align=center, fill=#1},
  hdr/.default=white,
  ll/.style={draw=gray!60, dashed, line width=0.5pt},
  arr/.style={-{Latex[length=2.3mm]}, line width=0.55pt},
  darr/.style={-{Latex[length=2.3mm]}, line width=0.55pt, dashed},
  actbar/.style={fill=blue!10, draw=blue!35, line width=0.4pt},
  compbox/.style={draw, rectangle, rounded corners=3pt,
                  minimum width=3.4cm, minimum height=0.65cm,
                  font=\small\ttfamily, fill=#1, align=center},
  compbox/.default=blue!8,
  ifbox/.style={draw, rectangle, rounded corners=3pt, dashed,
                minimum width=3.4cm, minimum height=0.65cm,
                font=\small\ttfamily\bfseries, fill=gray!10, align=center},
  clsbox/.style={draw, rectangle, rounded corners=2pt,
                 minimum width=2.6cm, minimum height=0.45cm,
                 font=\scriptsize\ttfamily, fill=blue!6, align=center},
  inh/.style={-{Triangle[open,length=2.5mm]}, line width=0.55pt}
}

% --------------- Comando testo codice ---------------
\newcommand{\code}[1]{\texttt{#1}}

% =============================================================================
\begin{document}
% =============================================================================

% ---- Pagina del titolo ----
\begin{titlepage}
\centering\vspace*{1cm}
\IfFileExists{logo_polimi.png}{%
  \includegraphics[width=0.35\textwidth]{logo_polimi.png}\\[1cm]}{}
{\large Scuola di Ingegneria Industriale e dell'Informazione}\\[0.5cm]
{\large Corso di Ingegneria del Software --- Prova Finale}\\[2cm]
{\huge\bfseries Mesos}\\[0.4cm]
{\Large Specifica del Protocollo di Rete}\\[0.4cm]
{\large Architettura Client-Server --- Socket TCP / Java RMI}\\[1.5cm]
\vfill
\begin{tabular}{rl}
\textbf{Docente:}  & Alessandro Margara \\[0.3cm]
\textbf{Studenti:} & Andrea Cacioppo   \\
                   & Andrea Carbone    \\
                   & Gabriele Buca     \\
                   & Francesco Botton  \\
\end{tabular}\\[1cm]
{\large Anno Accademico 2025/2026}
\vspace*{1cm}
\end{titlepage}

\tableofcontents
\newpage

% =============================================================================
\section{Panoramica del Protocollo}
% =============================================================================

Il protocollo di rete AM54 definisce le regole di comunicazione tra uno o più
client di gioco e un server centrale. Il sistema implementa un \textbf{doppio
strato di trasporto}: sia Socket TCP (object streams) sia Java RMI sono
supportati simultaneamente, condividendo la stessa gerarchia di messaggi e lo
stesso pattern Observer.

\subsection{Decisioni Progettuali}

\begin{description}[leftmargin=2.2cm, labelwidth=2cm]
  \item[\textbf{Dual Transport}]
    Lo stesso protocollo logico è disponibile su TCP (porta~12345) e
    RMI (porta~1099). La scelta è trasparente alla logica di gioco.

  \item[\textbf{Dispatch polimorfico}]
    I messaggi portano la propria logica di esecuzione: le \textit{Request}
    implementano \code{execute(Controller,\,String,\,ClientConnection)};
    gli \textit{Update} implementano \code{applyTo(ClientModel)}.
    Non sono necessari costrutti \code{switch}/\code{instanceof}.

  \item[\textbf{Observer Pattern}]
    La \code{VirtualView} lato server osserva il modello \code{Game} e
    traduce ogni evento in un \code{Update} trasmesso a tutti i client.

  \item[\textbf{Aggiornamenti differenziali}]
    La maggior parte degli update trasporta solo i dati cambiati.
    Lo \code{StateChangedUpdate} contiene lo stato completo ed è inviato
    ad ogni transizione di fase e alla riconnessione.
\end{description}

\subsection{Confronto dei Trasporti}

\begin{tabularx}{\textwidth}{@{}lXX@{}}
\toprule
\textbf{Aspetto} & \textbf{Socket TCP} & \textbf{RMI} \\
\midrule
Porta            & 12345 & 1099 \\
Serializzazione  & Java Object Streams & Protocollo RMI \\
Entry point      & \code{new Socket(host, 12345)} & \code{getRegistry(host,~1099)} \\
Consegna update  & Thread background legge ObjectInputStream & Server chiama stub del client \\
Heartbeat        & PingRequest / PongResponse & \code{ping()} / \code{pong()} remoti \\
Disconnessione   & \code{IOException} su socket & \code{RemoteException} \\
\bottomrule
\end{tabularx}

% =============================================================================
\section{Architettura}
% =============================================================================

\subsection{Struttura dei Package}

\begin{description}[leftmargin=0pt, font=\ttfamily\bfseries, itemsep=4pt]
  \item[network/client/]
    \begin{itemize}[label=\textendash, leftmargin=1.4cm, font=\ttfamily, itemsep=1pt]
      \item Client.java --- Singleton hub del client
      \item ServerConnection.java --- Astrazione trasporto (interfaccia)
      \item SocketServerConnection.java --- Implementazione TCP
      \item RMIServerConnection.java --- Implementazione RMI + target callback
      \item ServerHandler.java --- Thread lettura background TCP
      \item ClientModel.java --- Stato locale lato client
      \item VirtualServer.java --- Facciata azioni per la View
    \end{itemize}
  \item[network/server/]
    \begin{itemize}[label=\textendash, leftmargin=1.4cm, font=\ttfamily, itemsep=1pt]
      \item Server.java --- Entry point (ServerSocket TCP + registry RMI)
      \item ClientConnection.java --- Astrazione trasporto (interfaccia)
      \item SocketClientHandler.java --- Thread TCP per singolo client
      \item RMIClientHandler.java --- Handler RMI per singolo client
      \item VirtualView.java --- Osservatore Game → broadcast
      \item PersistenceManager.java --- Backup JSON per crash recovery
    \end{itemize}
  \item[network/rmi/]
    \begin{itemize}[label=\textendash, leftmargin=1.4cm, font=\ttfamily, itemsep=1pt]
      \item RMIGameServer.java --- Interfaccia remota d'ingresso
      \item RMIClientInterface.java --- Interfaccia callback client
      \item RMIClientHandlerInterface.java --- Interfaccia handler server
    \end{itemize}
  \item[network/messages/]
    \begin{itemize}[label=\textendash, leftmargin=1.4cm, font=\ttfamily, itemsep=1pt]
      \item ClientMessage.java --- Interfaccia marcatore radice (Serializable)
      \item Request.java --- Client $\to$ Server (\textit{sealed}, 9 sottotipi)
      \item Update.java --- Server $\to$ Client (\textit{sealed}, 16 sottotipi)
    \end{itemize}
\end{description}

\subsection{Architettura di Sistema}

\begin{figure}[H]
\centering
\begin{tikzpicture}[xscale=0.9]
  % CLIENT
  \node[compbox=clGreen]  (cli)  at (0,   0)    {Client.java};
  \node[ifbox]            (sc)   at (0,  -1.4)  {\small«interface»\\ServerConnection};
  \node[compbox=clGreen]  (ssc)  at (-2.5,-3.2) {SocketServer\\Connection};
  \node[compbox=clGreen]  (rsc)  at ( 2.5,-3.2) {RMIServer\\Connection};

  % SERVER
  \node[compbox=srvOrange] (srv)  at (10,  0)    {Server.java};
  \node[ifbox]             (cc)   at (10, -1.4)  {\small«interface»\\ClientConnection};
  \node[compbox=srvOrange] (sch)  at ( 7.5,-3.2) {SocketClient\\Handler};
  \node[compbox=srvOrange] (rmh)  at (12.5,-3.2) {RMIClient\\Handler};

  % MODELLO
  \node[compbox=mdlRed]   (ctrl) at (10,-4.8)   {Controller};
  \node[compbox=mdlRed]   (game) at (10,-6.2)   {Game (Model)};
  \node[compbox=mdlRed]   (vv)   at (10,-7.6)   {VirtualView};

  % Frecce client
  \draw[arr] (cli) -- (sc);
  \draw[arr] (sc.south) -- ++(0,-0.4) -| (ssc.north);
  \draw[arr] (sc.south) -- ++(0,-0.4) -| (rsc.north);

  % TCP/RMI
  \draw[arr] (ssc.east) -- node[above,font=\scriptsize]{TCP} (sch.west);
  \draw[arr] (rsc.east) -- node[above,font=\scriptsize]{RMI} (rmh.west);

  % Server chain
  \draw[arr] (sch.south) |- (ctrl.west);
  \draw[arr] (rmh.south) |- (ctrl.east);
  \draw[arr] (ctrl) -- (game);
  \draw[arr] (game) -- (vv);

  % Notifica VirtualView → ClientConnection
  \draw[arr] (vv.east) -- ++(1.2,0) -- ++(0,6.3) -| (cc.east);
  \node[right,font=\scriptsize\itshape] at (11.2,-4.0) {notifica};

  % ClientConnection → handler
  \draw[arr] (cc.south) -- ++(0,-0.4) -| (sch.north);
  \draw[arr] (cc.south) -- ++(0,-0.4) -| (rmh.north);

  \node[font=\small\bfseries, above] at (0,  0.5) {LATO CLIENT};
  \node[font=\small\bfseries, above] at (10, 0.5) {LATO SERVER};
\end{tikzpicture}
\caption{Architettura generale del sistema}
\end{figure}

% =============================================================================
\newpage
\section{Protocollo Messaggi}
% =============================================================================

\subsection{Gerarchia dei Messaggi}

\begin{figure}[H]
\centering
\begin{tikzpicture}

  % Radice
  \node[ifbox, minimum width=5cm] (root) at (5.5, 0)
      {«interface» ClientMessage\\{\scriptsize extends Serializable}};

  % Rami
  \node[ifbox, minimum width=3.5cm] (req) at (1.5, -2)
      {«sealed» Request\\{\scriptsize CLIENT $\to$ SERVER}};
  \node[ifbox, minimum width=3.5cm] (upd) at (9.5, -2)
      {«sealed» Update\\{\scriptsize SERVER $\to$ CLIENT}};

  \draw[inh] (req.north) -- ++(0,0.5) -| (root.south);
  \draw[inh] (upd.north) -- ++(0,0.5) -| (root.south);

  % Request subtypes
  \foreach \y/\lbl in {
    -3.5/JoinGameRequest,
    -4.2/SetupGameRequest,
    -4.9/ChooseTotemRequest,
    -5.6/PlaceTotemRequest,
    -6.3/PickCardRequest}
    \node[clsbox] at (-0.5, \y) {\lbl};

  \foreach \y/\lbl in {
    -3.5/EndTurnRequest,
    -4.2/GetLeaderboardRequest,
    -4.9/PingRequest,
    -5.6/PongResponse}
    \node[clsbox] at (3.5, \y) {\lbl};

  \draw[inh] (-0.5,-3.2) -- ++(0,0.35) -| (req.south);
  \draw[inh] ( 3.5,-3.2) -- ++(0,0.35) -| (req.south);

  % Update subtypes – colonna sinistra
  \foreach \y/\lbl in {
    -3.5/PlayerJoinedUpdate,
    -4.2/TotemColorUpdate,
    -4.9/TotemPlacedUpdate,
    -5.6/CardPickedUpdate,
    -6.3/TotemReturnedUpdate,
    -7.0/TurnChangedUpdate}
    \node[clsbox] at (7.0, \y) {\lbl};

  % Update subtypes – colonna destra
  \foreach \y/\lbl in {
    -3.5/StateChangedUpdate,
    -4.2/EndGameUpdate,
    -4.9/EventUpdate,
    -5.6/ErrorMessageUpdate,
    -6.3/PlayerDisconnectedUpdate,
    -7.0/PlayerReconnectedUpdate,
    -7.7/GameSuspendedUpdate,
    -8.4/GameResumedUpdate,
    -9.1/LeaderboardUpdate}
    \node[clsbox] at (12.0, \y) {\lbl};

  \draw[inh] (7.0, -3.2)  -- ++(0,0.35) -| (upd.south);
  \draw[inh] (12.0,-3.2)  -- ++(0,0.35) -| (upd.south);

\end{tikzpicture}
\caption{Gerarchia dei messaggi di rete}
\end{figure}

\subsection{Request: Client $\to$ Server}

\begin{small}
\begin{tabularx}{\textwidth}{@{}lp{4.2cm}X@{}}
\toprule
\textbf{Messaggio} & \textbf{Campi} & \textbf{Scopo} \\
\midrule
\code{JoinGameRequest}       & \code{String nickname}                & Connessione iniziale / rientro \\
\code{SetupGameRequest}      & ---                                   & Segnala prontezza ad iniziare \\
\code{ChooseTotemRequest}    & \code{TotemColors color}              & Scelta colore totem \\
\code{PlaceTotemRequest}     & \code{int offerPosition}              & Posiziona totem su offering tile \\
\code{PickCardRequest}       & \code{String row, int index}          & Preleva carta da una fila \\
\code{EndTurnRequest}        & ---                                   & Termina il turno corrente \\
\code{GetLeaderboardRequest} & \code{int numPlayers}                 & Richiede la classifica \\
\code{PingRequest}           & \code{long timestamp}                 & Heartbeat (è anche Update) \\
\code{PongResponse}          & \code{long timestamp}                 & Risposta all'heartbeat \\
\bottomrule
\end{tabularx}
\end{small}

\subsection{Update: Server $\to$ Client}

\begin{small}
\begin{tabularx}{\textwidth}{@{}lp{5cm}X@{}}
\toprule
\textbf{Messaggio} & \textbf{Campi principali} & \textbf{Scopo} \\
\midrule
\code{PlayerJoinedUpdate}       & \code{String nickname}               & Nuovo giocatore connesso \\
\code{TotemColorUpdate}         & \code{String nick, TotemColors c}    & Colore totem scelto \\
\code{TotemPlacedUpdate}        & \code{String nick, int pos}          & Totem posizionato \\
\code{CardPickedUpdate}         & nick, row, idx, food, magic, offering & Carta prelevata \\
\code{TotemReturnedUpdate}      & \code{String nick, int pos, food, aura} & Totem restituito \\
\code{TurnChangedUpdate}        & \code{String nickname}               & Cambio turno \\
\code{StateChangedUpdate}       & state, era, round, players, turnOrder, offering, 6 liste carte & Snapshot completo \\
\code{EndGameUpdate}            & \code{List winners, Map scores, int num} & Partita terminata \\
\code{EventUpdate}              & name, players, 6 liste               & Risoluzione evento \\
\code{ErrorMessageUpdate}       & \code{String message}                & Errore / mossa invalida \\
\code{PlayerDisconnectedUpdate} & \code{String nickname}               & Giocatore disconnesso \\
\code{PlayerReconnectedUpdate}  & \code{String nickname}               & Giocatore riconnesso \\
\code{GameSuspendedUpdate}      & ---                                  & Partita sospesa (1 attivo) \\
\code{GameResumedUpdate}        & ---                                  & Partita ripresa (\(\geq\)2) \\
\code{LeaderboardUpdate}        & \code{int pos, List entries, int num} & Risultato classifica \\
\code{PingRequest}              & \code{long timestamp}                & Heartbeat server\,$\to$\,client \\
\bottomrule
\end{tabularx}
\end{small}

% =============================================================================
\newpage
\section{Ciclo di Vita della Connessione}
% =============================================================================

\subsection{Sequenza di Avvio}

\begin{figure}[H]
\centering
\begin{tikzpicture}

  % --- headers ---
  \node[hdr=clGreen]   at (0,   0.4) {Client};
  \node[hdr=srvOrange] at (3.0, 0.4) {SocketCl\\Handler};
  \node[hdr=mdlRed]    at (6.0, 0.4) {Controller};
  \node[hdr=mdlRed]    at (9.0, 0.4) {Game};

  % --- lifelines ---
  \foreach \x in {0, 3.0, 6.0, 9.0}
    \draw[ll] (\x, 0) -- (\x, -8.5);

  % --- messages ---
  % TCP connect
  \draw[arr] (0,-1.0) -- node[above,font=\scriptsize\ttfamily]{TCP connect (Socket 12345)} (3.0,-1.0);

  % JoinGameRequest
  \draw[arr] (0,-2.0) -- node[above,font=\scriptsize\ttfamily]{JoinGameRequest(nickname)} (3.0,-2.0);

  % handleJoinGame
  \draw[arr] (3.0,-3.0) -- node[above,font=\scriptsize\ttfamily]{handleJoinGame(nick, conn)} (6.0,-3.0);
  \draw[actbar] (6.0,-3.0) rectangle (6.25,-5.5);

  % game.joinGame
  \draw[arr] (6.0,-3.7) -- node[above,font=\scriptsize\ttfamily]{game.joinGame(nick)} (9.0,-3.7);

  % new VirtualView
  \draw[arr] (6.0,-4.5) -- (6.0,-5.0);
  \node[right,font=\scriptsize\ttfamily] at (6.3,-4.75) {new VirtualView(conn)};

  % addObserver
  \draw[arr] (6.0,-5.5) -- node[above,font=\scriptsize\ttfamily]{game.addObserver(vv)} (9.0,-5.5);

  % startHeartbeat
  \draw[arr] (3.0,-6.8) -- (3.0,-7.3);
  \node[right,font=\scriptsize\ttfamily] at (3.1,-7.05) {connection.startHeartbeat()};

  % PlayerJoinedUpdate → client
  \draw[arr] (3.0,-8.0) -- node[above,font=\scriptsize\ttfamily]{PlayerJoinedUpdate} (0,-8.0);

\end{tikzpicture}
\caption{Sequenza di avvio: dal TCP connect alla notifica in lobby}
\end{figure}

\subsection{Ciclo di Vita di un'Azione}

Il diagramma traccia l'azione \textit{PlaceTotem} dal click sulla UI
all'aggiornamento di tutti i client connessi.

\begin{figure}[H]
\centering
\begin{tikzpicture}

  \node[hdr=clGreen]   at (0,   0.4) {View / UI};
  \node[hdr=clGreen]   at (2.6, 0.4) {Virtual\\Server};
  \node[hdr=clGreen]   at (5.2, 0.4) {ServerConn.};
  \node[hdr=srvOrange] at (7.8, 0.4) {SocketCl\\Handler};
  \node[hdr=mdlRed]    at (10.4,0.4) {Controller};
  \node[hdr=mdlRed]    at (13.0,0.4) {Game};

  \foreach \x in {0, 2.6, 5.2, 7.8, 10.4, 13.0}
    \draw[ll] (\x, 0) -- (\x, -11.5);

  % 1. placeTotem
  \draw[arr] (0,-1.1) -- node[above,font=\scriptsize\ttfamily]{placeTotem(3)} (2.6,-1.1);

  % 2. getConnection via Client singleton
  \draw[arr] (2.6,-2.0) -- node[above,font=\scriptsize\ttfamily]{connection() $\to$ Client.getInstance()} (2.6,-2.5);

  % 3. sendObject
  \draw[arr] (2.6,-3.1) -- node[above,font=\scriptsize\ttfamily]{sendObject(PlaceTotemRequest(3))} (5.2,-3.1);

  % 4. TCP stream
  \draw[arr] (5.2,-4.1) -- node[above,font=\scriptsize\ttfamily]{ObjectOutputStream.writeObject} (7.8,-4.1);

  % 5. readObject
  \draw[arr] (7.8,-5.0) -- (7.8,-5.5);
  \node[right,font=\scriptsize\ttfamily] at (7.9,-5.25) {readObject()};

  % 6. request.execute
  \draw[arr] (7.8,-6.2) -- node[above,font=\scriptsize\ttfamily]{request.execute(ctrl, nick, this)} (10.4,-6.2);
  \draw[actbar] (10.4,-6.2) rectangle (10.65,-7.3);

  % 7. game.placeTotem
  \draw[arr] (10.4,-7.0) -- node[above,font=\scriptsize\ttfamily]{game.placeTotem(nick, 3)} (13.0,-7.0);
  \draw[actbar] (13.0,-7.0) rectangle (13.25,-8.2);

  % 8. notifyObservers
  \draw[arr] (13.0,-7.9) -- (13.0,-8.3);
  \node[right,font=\scriptsize\ttfamily] at (13.1,-8.1) {notifyObservers()};

  % 9. VirtualView → ClientConnection
  \draw[arr] (10.4,-9.0) -- node[above,font=\scriptsize\ttfamily]{connection.notifyTotemPlaced(nick, pos)} (7.8,-9.0);

  % 10. ObjectOutputStream → client
  \draw[arr] (7.8,-9.9) -- node[above,font=\scriptsize\ttfamily]{TotemPlacedUpdate} (5.2,-9.9);

  % 11. applyTo
  \draw[arr] (5.2,-10.8) -- node[above,font=\scriptsize\ttfamily]{update.applyTo(ClientModel)} (2.6,-10.8);

\end{tikzpicture}
\caption{Ciclo di vita completo di un'azione (PlaceTotemRequest)}
\end{figure}

% =============================================================================
\newpage
\section{Heartbeat e Disconnessione}
% =============================================================================

\subsection{Meccanismo Ping-Pong}

Il server usa uno \code{ScheduledExecutorService} per inviare un ping ogni
\textbf{5\,000\,ms}. Se entro \textbf{15\,000\,ms} non arriva risposta,
la connessione è dichiarata persa.

\begin{figure}[H]
\centering
\begin{tikzpicture}

  \node[hdr=clGreen]   at (0,   0.4) {Client};
  \node[hdr=srvOrange] at (5.0, 0.4) {SocketCl\\Handler};
  \node[hdr=mdlRed]    at (10.0,0.4) {Controller};

  \foreach \x in {0, 5.0, 10.0}
    \draw[ll] (\x, 0) -- (\x, -9.5);

  % startHeartbeat
  \draw[arr] (5.0,-0.9) -- (5.0,-1.4);
  \node[right,font=\scriptsize\ttfamily] at (5.1,-1.15)
      {startHeartbeat() scheduleAtFixedRate(5\,000 ms)};

  % Primo ciclo
  \draw[arr] (5.0,-2.4) -- node[above,font=\scriptsize\ttfamily]{PingRequest} (0,-2.4);
  \draw[arr] (0,-3.1) -- node[above,font=\scriptsize\ttfamily]{PongResponse} (5.0,-3.1);
  \draw[arr] (5.0,-3.6) -- (5.0,-4.0);
  \node[right,font=\scriptsize\ttfamily] at (5.1,-3.8) {lastPongTime = now()};

  \node[font=\scriptsize\itshape] at (2.5,-4.7) {\ldots\ ogni 5\,s\ \ldots};

  % Secondo ciclo
  \draw[arr] (5.0,-5.6) -- node[above,font=\scriptsize\ttfamily]{PingRequest} (0,-5.6);
  \draw[arr] (0,-6.3) -- node[above,font=\scriptsize\ttfamily]{PongResponse} (5.0,-6.3);

  % Timeout
  \draw[arr] (5.0,-7.3) -- (5.0,-7.8);
  \node[right,font=\scriptsize\ttfamily,text=red!60!black]
      at (5.1,-7.55) {timeSincePong $>$ 15\,000 ms $\to$ timeout};

  \draw[arr] (5.0,-8.6) -- node[above,font=\scriptsize\ttfamily]{removeClient(nick)} (10.0,-8.6);

\end{tikzpicture}
\caption{Heartbeat ping-pong e rilevamento timeout}
\end{figure}

\subsection{Flusso di Disconnessione}

\begin{figure}[H]
\centering
\begin{tikzpicture}

  \node[hdr=srvOrange] at (0,   0.4) {Client\\Handler};
  \node[hdr=mdlRed]    at (3.2, 0.4) {Controller};
  \node[hdr=sysGray]   at (6.4, 0.4) {Disconn.\\Set};
  \node[hdr=mdlRed]    at (9.6, 0.4) {VirtualView\\(altri)};
  \node[hdr=clGreen]   at (12.8,0.4) {Altri\\Client};

  \foreach \x in {0, 3.2, 6.4, 9.6, 12.8}
    \draw[ll] (\x, 0) -- (\x, -10.5);

  % removeClient
  \draw[arr] (0,-1.2) -- node[above,font=\scriptsize\ttfamily]{removeClient(nick)} (3.2,-1.2);
  \draw[actbar] (3.2,-1.2) rectangle (3.45,-3.3);

  % add to set
  \draw[arr] (3.2,-1.9) -- node[above,font=\scriptsize\ttfamily]{disconnectedPlayers.add(nick)} (6.4,-1.9);

  % remove observer
  \draw[arr] (3.2,-2.7) -- (3.2,-3.3);
  \node[right,font=\scriptsize\ttfamily] at (3.5,-3.0) {game.removeObserver(vv)};

  % broadcast
  \draw[arr] (3.2,-4.2) -- node[above,font=\scriptsize\ttfamily]{broadcastPlayerDisconnected(nick)} (9.6,-4.2);
  \draw[actbar] (9.6,-4.2) rectangle (9.85,-4.8);
  \draw[arr] (9.6,-4.7) -- node[above,font=\scriptsize\ttfamily]{PlayerDisconnectedUpdate} (12.8,-4.7);

  % autoHandleDisconnection
  \draw[arr] (3.2,-5.8) -- (3.2,-6.3);
  \node[right,font=\scriptsize\ttfamily] at (3.5,-6.05) {autoHandleDisconnection()};

  % skipTurn
  \draw[arr] (3.2,-7.0) -- (3.2,-7.5);
  \node[right,font=\scriptsize\ttfamily] at (3.5,-7.25) {skipTurn(game, nick) se è il turno del giocatore};

  % checkSuspension
  \draw[arr] (3.2,-8.2) -- (3.2,-8.7);
  \node[right,font=\scriptsize\ttfamily] at (3.5,-8.45) {checkSuspension()};

  % se 1 solo attivo
  \draw[arr] (3.2,-9.5) -- node[above,font=\scriptsize\ttfamily]{broadcastGameSuspended()} (9.6,-9.5);
  \draw[actbar] (9.6,-9.5) rectangle (9.85,-10.1);
  \draw[arr] (9.6,-10.0) -- node[above,font=\scriptsize\ttfamily]{GameSuspendedUpdate} (12.8,-10.0);

\end{tikzpicture}
\caption{Cascata di disconnessione durante una partita attiva}
\end{figure}

\subsection{Sospensione e Vittoria per Abbandono}

Con un solo giocatore attivo il server avvia un countdown di 60\,s.
Allo scadere il giocatore rimasto vince per abbandono.

\begin{figure}[H]
\centering
\begin{tikzpicture}

  \node[hdr=mdlRed]    at (0,   0.4) {Controller};
  \node[hdr=mdlRed]    at (3.2, 0.4) {Game};
  \node[hdr=sysGray]   at (6.4, 0.4) {Scheduler};
  \node[hdr=sysGray]   at (9.6, 0.4) {Database};
  \node[hdr=srvOrange] at (12.8,0.4) {Persistence};

  \foreach \x in {0, 3.2, 6.4, 9.6, 12.8}
    \draw[ll] (\x, 0) -- (\x, -9.5);

  % schedule
  \draw[arr] (0,-1.2) -- node[above,font=\scriptsize\ttfamily]{schedule(60\,s)} (6.4,-1.2);

  % timeout
  \draw[arr] (6.4,-2.5) -- node[above,font=\scriptsize\ttfamily]{timeout $\to$ suspendGame()} (0,-2.5);
  \draw[actbar] (0,-2.5) rectangle (0.25,-7.8);

  % setState
  \draw[arr] (0,-3.3) -- node[above,font=\scriptsize\ttfamily]{setState(EndGameState)} (3.2,-3.3);

  % captureFinalResult
  \draw[arr] (0,-4.2) -- (0,-4.7);
  \node[right,font=\scriptsize\ttfamily] at (0.3,-4.45) {captureFinalResult(winner)};

  % save DB
  \draw[arr] (0,-5.5) -- node[above,font=\scriptsize\ttfamily]{saveMatchResult(nick, score, n)} (9.6,-5.5);

  % clearBackup
  \draw[arr] (0,-6.5) -- node[above,font=\scriptsize\ttfamily]{persistenceManager.clearBackup()} (12.8,-6.5);

  % notifyEndGame
  \draw[arr] (0,-7.5) -- node[above,font=\scriptsize\ttfamily]{game.notifyEndGame(winner)} (3.2,-7.5);

  % EndGameUpdate → client
  \draw[arr] (3.2,-8.6) -- node[above,font=\scriptsize\ttfamily]{EndGameUpdate (via VirtualView)} (0,-8.6);

\end{tikzpicture}
\caption{Timeout da sospensione e vittoria per abbandono}
\end{figure}

% =============================================================================
\newpage
\section{Protocollo di Riconnessione}
% =============================================================================

\subsection{Flusso di Riconnessione Server-Side}

\begin{figure}[H]
\centering
\begin{tikzpicture}

  \node[hdr=clGreen]   at (0,   0.4) {Client};
  \node[hdr=srvOrange] at (2.8, 0.4) {Client\\Handler};
  \node[hdr=mdlRed]    at (5.6, 0.4) {Controller};
  \node[hdr=sysGray]   at (8.4, 0.4) {Disconn.\\Set};
  \node[hdr=mdlRed]    at (11.2,0.4) {Game};

  \foreach \x in {0, 2.8, 5.6, 8.4, 11.2}
    \draw[ll] (\x, 0) -- (\x, -12.0);

  % JoinGameRequest
  \draw[arr] (0,-1.1) -- node[above,font=\scriptsize\ttfamily]{JoinGameRequest(nick)} (2.8,-1.1);
  \draw[arr] (2.8,-1.8) -- node[above,font=\scriptsize\ttfamily]{handleJoinGame(nick, conn)} (5.6,-1.8);
  \draw[actbar] (5.6,-1.8) rectangle (5.85,-4.5);

  % check disconnected
  \draw[arr] (5.6,-2.5) -- node[above,font=\scriptsize\ttfamily]{contains(nick)?} (8.4,-2.5);
  \draw[darr] (8.4,-3.1) -- node[above,font=\scriptsize\ttfamily]{YES $\to$ handleReconnect()} (5.6,-3.1);

  % remove from set
  \draw[arr] (5.6,-3.9) -- node[above,font=\scriptsize\ttfamily]{disconnectedPlayers.remove(nick)} (8.4,-3.9);

  % swap VirtualView
  \draw[arr] (5.6,-5.2) -- node[above,font=\scriptsize\ttfamily]{game.removeObserver(oldVV)} (11.2,-5.2);
  \draw[arr] (5.6,-6.0) -- node[above,font=\scriptsize\ttfamily]{new VirtualView(conn, nick)} (11.2,-6.0);
  \draw[arr] (5.6,-6.8) -- node[above,font=\scriptsize\ttfamily]{game.addObserver(newVV)} (11.2,-6.8);

  % snapshot
  \draw[arr] (5.6,-7.8) -- node[above,font=\scriptsize\ttfamily]{notifyAllClientsStateChanged()} (11.2,-7.8);
  \draw[arr] (5.6,-8.6) -- node[above,font=\scriptsize\ttfamily]{StateChangedUpdate (snapshot completo)} (2.8,-8.6);
  \draw[arr] (2.8,-9.2) -- node[above,font=\scriptsize\ttfamily]{ClientModel.replaceState()} (0,-9.2);

  % onTurnChanged
  \draw[arr] (5.6,-10.0) -- node[above,font=\scriptsize\ttfamily]{onTurnChanged(currentPlayer)} (11.2,-10.0);

  % EndGameUpdate se terminata
  \draw[arr] (5.6,-10.8) -- node[above,font=\scriptsize\ttfamily]{EndGameUpdate (se partita terminata)} (2.8,-10.8);

  % broadcastPlayerReconnected
  \draw[arr] (5.6,-11.6) -- node[above,font=\scriptsize\ttfamily]{broadcastPlayerReconnected(nick)} (11.2,-11.6);

\end{tikzpicture}
\caption{Riconnessione: il server ripristina lo stato tramite snapshot completo}
\end{figure}

\subsection{Auto-Riconnessione Lato Client}

\begin{figure}[H]
\centering
\begin{tikzpicture}

  \node[hdr=clGreen]   at (0,   0.4) {Client};
  \node[hdr=clGreen]   at (3.5, 0.4) {Reconnect\\Thread};
  \node[hdr=srvOrange] at (7.0, 0.4) {SocketCl\\Handler};

  \foreach \x in {0, 3.5, 7.0}
    \draw[ll] (\x, 0) -- (\x, -8.5);

  % Exception
  \draw[arr] (7.0,-1.1) -- node[above,font=\scriptsize\ttfamily]{IOException / RemoteException} (0,-1.1);

  % startAutoReconnect
  \draw[arr] (0,-2.0) -- node[above,font=\scriptsize\ttfamily]{startAutoReconnect()} (3.5,-2.0);

  % loop
  \draw[dashed, rounded corners=3pt] (3.1,-2.7) rectangle (7.4,-8.0);
  \node[anchor=north west,font=\scriptsize\bfseries] at (3.2,-2.8)
      {loop [ogni 3\,s]};

  \draw[arr] (3.5,-3.7) -- (3.5,-4.2);
  \node[right,font=\scriptsize\ttfamily] at (3.6,-3.95) {new Socket(ip, 12345)};

  \draw[arr] (3.5,-5.0) -- node[above,font=\scriptsize\ttfamily]{TCP connect} (7.0,-5.0);
  \draw[arr] (3.5,-5.8) -- node[above,font=\scriptsize\ttfamily]{JoinGameRequest(nickname)} (7.0,-5.8);
  \draw[darr] (7.0,-6.6) -- node[above,font=\scriptsize\ttfamily]{StateChangedUpdate (success)} (3.5,-6.6);

  % aggiorna Client singleton
  \draw[arr] (3.5,-7.5) -- node[above,font=\scriptsize\ttfamily]{Client.setConnection(newConn)} (0,-7.5);

\end{tikzpicture}
\caption{Auto-riconnessione client: retry con nuova connessione ogni 3\,s}
\end{figure}

% =============================================================================
\newpage
\section{Recupero da Crash del Server}
% =============================================================================

\subsection{Backup JSON e Ripristino}

Il server persiste l'intero oggetto \code{Game} su \code{backup.json} dopo
ogni azione che altera lo stato. Il salvataggio è atomico tramite
\code{Files.move(ATOMIC\_MOVE)}: il file temporaneo viene completamente scritto
prima di sostituire l'originale, garantendo che un crash durante la scrittura
non corrompa il backup precedente.

Il backup usa Gson con \code{RuntimeTypeAdapterFactory} per la serializzazione
polimorfica di \code{GameState}, \code{IEffect} e \code{TribeCard}.
Il metodo \code{restoreReferences()} ri-collega i puntatori tra oggetti
che Gson ricostruisce come copie separate.

\begin{figure}[H]
\centering
\begin{tikzpicture}

  \node[hdr=mdlRed]    at (0,   0.4) {Controller};
  \node[hdr=sysGray]   at (3.5, 0.4) {Persistence\\Manager};
  \node[hdr=sysGray]   at (7.0, 0.4) {File System};
  \node[hdr=srvOrange] at (10.5,0.4) {Server\\(riavvio)};

  \foreach \x in {0, 3.5, 7.0, 10.5}
    \draw[ll] (\x, 0) -- (\x, -8.0);

  % saveState
  \draw[arr] (0,-1.1) -- node[above,font=\scriptsize\ttfamily]{saveState()} (3.5,-1.1);
  \draw[actbar] (3.5,-1.1) rectangle (3.75,-3.2);

  \draw[arr] (3.5,-1.8) -- node[above,font=\scriptsize\ttfamily]{gson.toJson(game) $\to$ backup.json.tmp} (7.0,-1.8);
  \draw[arr] (3.5,-2.7) -- node[above,font=\scriptsize\ttfamily]{Files.move(tmp, backup, ATOMIC\_MOVE)} (7.0,-2.7);

  % CRASH
  \draw[line width=1.2pt, red!55, decorate,
        decoration={snake, amplitude=1mm, segment length=5mm}]
    (-0.3,-4.0) -- (10.8,-4.0);
  \node[above,font=\small\bfseries,text=red!65!black]
      at (5.25,-3.95) {SERVER CRASH};

  % Riavvio
  \draw[arr] (10.5,-5.2) -- node[above,font=\scriptsize\ttfamily]{new Server() $\to$ loadGame()} (3.5,-5.2);
  \draw[actbar] (3.5,-5.2) rectangle (3.75,-7.2);

  \draw[arr] (3.5,-5.9) -- node[above,font=\scriptsize\ttfamily]{gson.fromJson(reader, Game.class)} (7.0,-5.9);
  \draw[arr] (3.5,-6.8) -- (3.5,-7.3);
  \node[right,font=\scriptsize\ttfamily] at (3.8,-7.05) {game.restoreReferences()};

\end{tikzpicture}
\caption{Backup atomico e recupero dopo crash del server}
\end{figure}

% =============================================================================
\newpage
\section{Interfacce di Connessione Lato Server}
% =============================================================================

\subsection{Gerarchia ClientConnection}

\begin{figure}[H]
\centering
\begin{tikzpicture}

  \node[ifbox, minimum width=5cm] (iface) at (5.5, 0)
      {«interface» ClientConnection};

  \node[compbox=srvOrange, minimum width=4.5cm] (sch)  at (2.0,-2.0)
      {SocketClientHandler};
  \node[compbox=srvOrange, minimum width=4.5cm] (rmih) at (9.0,-2.0)
      {RMIClientHandler};

  \draw[inh] (sch.north)  -- ++(0,0.5) -| (iface.south);
  \draw[inh] (rmih.north) -- ++(0,0.5) -| (iface.south);

  \node[anchor=north, font=\small, align=left] at (2.0,-2.8)
      {\textendash\ Scrive Update su ObjectOutputStream\\
       \textendash\ Thread TCP dedicato per client\\
       \textendash\ Heartbeat via PingRequest/PongResponse};

  \node[anchor=north, font=\small, align=left] at (9.0,-2.8)
      {\textendash\ Chiama stub RMIClientInterface\\
       \textendash\ Handler RMI per singolo client\\
       \textendash\ Heartbeat via metodi remoti ping()/pong()};

\end{tikzpicture}
\caption{Gerarchia di ClientConnection lato server}
\end{figure}

\subsection{API ClientConnection}

\begin{small}
\begin{tabularx}{\textwidth}{@{}p{6.5cm}X@{}}
\toprule
\textbf{Metodo} & \textbf{Descrizione} \\
\midrule
\code{notifyPlayerJoined(nickname)}               & Nuovo giocatore entrato in lobby \\
\code{notifyTotemColor(nick, color)}               & Colore totem scelto \\
\code{notifyTotemPlaced(nick, pos)}                & Totem posizionato su offering tile \\
\code{notifyCardPicked(nick, row, idx, ...)}        & Carta prelevata dal mercato \\
\code{notifyTotemReturned(nick, pos, food, aura)}  & Totem restituito a fine pesca \\
\code{notifyEventResolved(name, stats, ...)}        & Evento risolto \\
\code{notifyTurnChanged(nickname)}                 & Cambio turno \\
\code{notifyStateChanged(state, era, round, ...)}   & Snapshot completo di stato \\
\code{notifyEndGame(winners, scores, num)}          & Partita terminata \\
\code{notifyError(message)}                        & Errore o mossa non valida \\
\code{notifyPlayerDisconnected(nickname)}          & Giocatore disconnesso \\
\code{notifyPlayerReconnected(nickname)}           & Giocatore riconnesso \\
\code{notifyGameSuspended()}                       & Partita sospesa (1 attivo) \\
\code{notifyGameResumed()}                         & Partita ripresa (\(\geq 2\) attivi) \\
\code{sendPing()}                                  & Invia heartbeat al client \\
\code{sendLeaderboardUpdate(update)}               & Risultato classifica \\
\code{disconnect()}                                & Chiude la connessione \\
\code{startHeartbeat()}                            & Avvia il ping scheduler \\
\bottomrule
\end{tabularx}
\end{small}

% =============================================================================
\newpage
\section{Interfacce di Connessione Lato Client}
% =============================================================================

\subsection{ServerConnection}

\begin{small}
\begin{tabularx}{\textwidth}{@{}p{5.5cm}X@{}}
\toprule
\textbf{Metodo} & \textbf{Descrizione} \\
\midrule
\code{joinGame(nickname)}            & Connessione iniziale / rientro dopo crash \\
\code{setupGame()}                   & Segnala prontezza ad iniziare \\
\code{chooseTotem(color)}            & Sceglie il colore del totem \\
\code{placeTotem(position)}          & Posiziona il totem su una offering tile \\
\code{pickCard(row, index)}          & Preleva una carta dalla fila indicata \\
\code{endTurn()}                     & Termina il proprio turno \\
\code{requestLeaderboard(numPlayers)} & Richiede la classifica filtrata \\
\code{disconnect()}                  & Chiude la connessione \\
\bottomrule
\end{tabularx}
\end{small}

\subsection{Implementazioni}

\begin{figure}[H]
\centering
\begin{tikzpicture}

  \node[ifbox, minimum width=5cm] (sc) at (5.5, 0) {«interface» ServerConnection};

  \node[compbox=clGreen, minimum width=4.5cm] (ssc) at (2.0,-2.0)
      {SocketServerConnection};
  \node[compbox=clGreen, minimum width=4.5cm] (rsc) at (9.0,-2.0)
      {RMIServerConnection};

  \draw[inh] (ssc.north) -- ++(0,0.5) -| (sc.south);
  \draw[inh] (rsc.north) -- ++(0,0.5) -| (sc.south);

  \node[ifbox, minimum width=4.5cm] (rmici) at (9.0,-4.0)
      {«interface» RMIClientInterface};
  \draw[inh] (rsc.south) -- (rmici.north);
  \node[right,font=\scriptsize\itshape] at (9.2,-3.0) {target callback RMI};

  \node[anchor=north,font=\small,align=left] at (2.0,-2.8)
      {\textendash\ Scrive Request su ObjectOutputStream\\
       \textendash\ Avvia thread ServerHandler};

\end{tikzpicture}
\caption{RMIServerConnection implementa anche RMIClientInterface per ricevere le callback}
\end{figure}

\subsection{VirtualServer --- Facciata per la View}

\code{VirtualServer} recupera la connessione corrente dal singleton \code{Client}
ad ogni chiamata (\code{Client.getInstance().getConnection()}), non al momento
della costruzione. Dopo un'auto-riconnessione che sostituisce la
\code{ServerConnection}, tutte le azioni future vengono automaticamente
instradate sulla nuova connessione senza che la View debba essere aggiornata.

\begin{small}
\begin{tabularx}{\textwidth}{@{}p{5cm}X@{}}
\toprule
\textbf{Metodo} & \textbf{Comportamento} \\
\midrule
\code{joinGame()}            & \code{connection()} $\to$ \code{joinGame(nickname)} \\
\code{setupGame()}           & \code{connection()} $\to$ \code{setupGame()} \\
\code{chooseTotem(color)}    & \code{connection()} $\to$ \code{chooseTotem(color)} \\
\code{placeTotem(pos)}       & \code{connection()} $\to$ \code{placeTotem(pos)} \\
\code{pickCard(row, idx)}    & \code{connection()} $\to$ \code{pickCard(row, idx)} \\
\code{endTurn()}             & \code{connection()} $\to$ \code{endTurn()} \\
\code{requestLeaderboard(n)} & \code{connection()} $\to$ \code{requestLeaderboard(n)} \\
\bottomrule
\end{tabularx}
\end{small}

% =============================================================================
\newpage
\section{Protocollo TCP Socket}
% =============================================================================

\subsection{Layout degli Stream}

Ogni connessione TCP utilizza due \code{ObjectOutputStream} indipendenti su
un singolo socket: uno per le Request (Client$\,\to\,$Server) e uno per gli
Update (Server$\,\to\,$Client).

\begin{figure}[H]
\centering
\begin{tikzpicture}[scale=0.9]

  \node[compbox=clGreen, minimum width=3.0cm] (cs) at (0,  0.8) {Client invia\\Request};
  \node[compbox=clGreen, minimum width=3.0cm] (cr) at (0, -0.8) {Client riceve\\Update};

  \draw[draw=gray, fill=gray!8, rounded corners=2pt] (4.2,-1.4) rectangle (6.8,1.4);
  \node[font=\small\bfseries] at (5.5, 0) {Socket};
  \node[font=\scriptsize\ttfamily, rotate=90] at (4.7, 0) {Request stream};
  \node[font=\scriptsize\ttfamily, rotate=90] at (6.3, 0) {Update stream};

  \node[compbox=srvOrange, minimum width=3.0cm] (ss) at (11, 0.8) {Server invia\\Update};
  \node[compbox=srvOrange, minimum width=3.0cm] (sr) at (11,-0.8) {Server riceve\\Request};

  \draw[arr] (cs.east) -- node[above,font=\scriptsize\ttfamily]{writeObject(Request)} (4.2, 0.8);
  \draw[arr] (6.8, 0.8) -- node[above,font=\scriptsize\ttfamily]{writeObject(Update)} (ss.west);
  \draw[arr] (6.8,-0.8) -- node[below,font=\scriptsize\ttfamily]{readObject()} (sr.west);
  \draw[arr] (4.2,-0.8) -- node[below,font=\scriptsize\ttfamily]{readObject()} (cr.east);

\end{tikzpicture}
\caption{Doppio stream su singolo socket TCP}
\end{figure}

\subsection{Dispatch Lato Server e Lato Client}

\textbf{Server} (\code{SocketClientHandler.run()}):
\begin{enumerate}[leftmargin=*, label=\textbf{\arabic*.}]
  \item \textbf{Intercettazione precoce} --- \code{PongResponse} aggiorna
    \code{lastPongTime}; \code{JoinGameRequest} cattura il nickname.
  \item \textbf{Dispatch polimorfico} --- tutti gli altri sottotipi invocano
    \code{request.execute(controller, nick, this)}.
\end{enumerate}

\textbf{Client} (\code{ServerHandler.handleUpdate()}):
\begin{enumerate}[leftmargin=*, label=\textbf{\arabic*.}]
  \item \textbf{Intercettazione precoce} --- \code{PingRequest} è risposto
    automaticamente con \code{PongResponse}, in modo trasparente alla UI.
  \item \textbf{Dispatch polimorfico} --- tutti gli altri sottotipi invocano
    \code{message.applyTo(clientModel)}.
\end{enumerate}

% =============================================================================
\newpage
\section{Interfacce Remote RMI}
% =============================================================================

\subsection{Relazioni tra le Interfacce}

\begin{figure}[H]
\centering
\begin{tikzpicture}

  \node[ifbox, minimum width=5cm] (gs) at (5.5,0)
      {«interface» RMIGameServer\\{\scriptsize punto d'ingresso nel registry RMI}};

  \node[ifbox, minimum width=4.8cm] (chi) at (1.5,-2.5)
      {«interface»\\RMIClientHandlerInterface};
  \node[ifbox, minimum width=4.8cm] (ci)  at (9.5,-2.5)
      {«interface»\\RMIClientInterface};

  \node[compbox=srvOrange, minimum width=4.5cm] (impl) at (1.5,-4.5)
      {RMIClientHandler};

  \draw[inh] (chi.north) -- ++(0,0.4) -| (gs.south);
  \node[font=\scriptsize\itshape] at (3.5,-1.0) {connect() ritorna uno stub handler};

  \draw[inh] (impl.north) -- (chi.south);
  \draw[darr] (impl.east) -- node[above,font=\scriptsize\itshape]{implementa anche} (ci.south |- impl.north) -- (ci.south);

  \node[draw=blue!25, fill=noteBlue, rounded corners=3pt,
        inner sep=5pt, font=\small, align=left] at (9.5,-5.5)
    {\textbf{Callback pattern}\\
     Il server chiama metodi di RMIClientInterface\\
     direttamente sul client per consegnare gli update.\\
     Equivalente RMI di writeObject(Update) su TCP.};

\end{tikzpicture}
\caption{Relazioni tra le tre interfacce remote RMI}
\end{figure}

\subsection{RMIGameServer --- Punto di Ingresso}

\begin{small}
\begin{tabularx}{\textwidth}{@{}p{7cm}X@{}}
\toprule
\textbf{Metodo} & \textbf{Descrizione} \\
\midrule
\code{connect(RMIClientInterface clientStub)} &
  Registra il client, restituisce uno stub \code{RMIClientHandlerInterface} \\
\bottomrule
\end{tabularx}
\end{small}

\subsection{RMIClientHandlerInterface --- Azioni di Gioco}

\begin{small}
\begin{tabularx}{\textwidth}{@{}p{6cm}X@{}}
\toprule
\textbf{Metodo} & \textbf{Descrizione} \\
\midrule
\code{joinGameRequest(nickname)}         & Connessione / rientro \\
\code{setupGameRequest()}                & Prontezza ad iniziare \\
\code{chooseTotemColorRequest(color)}    & Scelta colore totem \\
\code{placeTotemRequest(offerPosition)}  & Posizionamento totem \\
\code{pickCardRequest(row, cardIndex)}   & Pesca carta \\
\code{endTurnRequest()}                  & Fine turno \\
\code{getLeaderboardRequest(numPlayers)} & Richiesta classifica \\
\code{pong()}                            & Risposta heartbeat \\
\code{disconnect()}                      & Disconnessione esplicita \\
\bottomrule
\end{tabularx}
\end{small}

\subsection{RMIClientInterface --- Callback Client}

\begin{small}
\begin{tabularx}{\textwidth}{@{}p{7cm}X@{}}
\toprule
\textbf{Metodo} & \textbf{Descrizione} \\
\midrule
\code{onPlayerJoined(nickname)}          & Nuovo giocatore in lobby \\
\code{onTotemColorChosen(nick, color)}   & Colore totem scelto \\
\code{onTotemPlaced(nick, pos)}          & Totem posizionato \\
\code{onTotemReturned(nick, food, aura)} & Totem restituito \\
\code{onCardPicked(nick, row, idx, food, magic, offering)} & Carta prelevata \\
\code{onTurnChanged(nickname)}           & Cambio turno \\
\code{onStateChanged(state, era, round, players, ...)} & Snapshot completo \\
\code{onEndGame(EndGameUpdate)}          & Partita terminata \\
\code{onErrorMessage(message)}           & Errore \\
\code{onLeaderboardUpdate(update)}       & Risultato classifica \\
\code{onPlayerDisconnected(nickname)}    & Giocatore disconnesso \\
\code{onPlayerReconnected(nickname)}     & Giocatore riconnesso \\
\code{onGameSuspended()}                 & Partita sospesa \\
\code{onGameResumed()}                   & Partita ripresa \\
\code{ping()}                            & Heartbeat server $\to$ client \\
\bottomrule
\end{tabularx}
\end{small}

% =============================================================================
\newpage
\section{Thread Safety}
% =============================================================================

\begin{description}[leftmargin=2.2cm, labelwidth=2cm, itemsep=6pt]
  \item[\textbf{Controller}]
    Tutte le azioni di gioco tramite \code{executeCommand} sono sincronizzate
    sul monitor \code{game.getPlayers()}, impedendo mutazioni concorrenti da
    thread client multipli.

  \item[\textbf{ClientModel}]
    Tutti i metodi di aggiornamento sono \code{synchronized}. Il thread
    \code{ServerHandler} (TCP) o il thread di callback RMI aggiornano il
    modello in modo sicuro rispetto al thread UI JavaFX.

  \item[\textbf{Socket I/O}]
    Il metodo \code{writeToStream} in \code{SocketClientHandler} è
    \code{synchronized} per prevenire interleaving di scritture sullo stesso
    \code{ObjectOutputStream} da parte di observer multipli.

  \item[\textbf{Heartbeat}]
    Il campo \code{lastPongTime} è \code{volatile long} per garantire
    visibilità immediata tra il thread di callback e il ping scheduler.

  \item[\textbf{Stato finale}]
    I campi \code{finalWinners}, \code{finalScores} e \code{finalNumPlayers}
    del \code{Controller} sono \code{volatile} per rendere visibile il
    risultato catturato a fine partita al thread che gestisce la riconnessione.
\end{description}

% =============================================================================
\section{Riferimenti ai File Sorgente}
% =============================================================================

\begin{small}
\begin{tabularx}{\textwidth}{@{}p{7.5cm}X@{}}
\toprule
\textbf{File} & \textbf{Ruolo} \\
\midrule
\code{network/server/Server.java}               & Entry point: ServerSocket TCP + registry RMI \\
\code{network/server/SocketClientHandler.java}  & Thread TCP per singolo client \\
\code{network/server/RMIServerImpl.java}        & Implementazione server RMI \\
\code{network/server/RMIClientHandler.java}     & Handler RMI per singolo client \\
\code{network/server/VirtualView.java}          & Osservatore Game $\to$ broadcast di rete \\
\code{network/server/ClientConnection.java}     & Interfaccia astrazione trasporto server \\
\code{network/server/PersistenceManager.java}   & Backup JSON per crash recovery \\
\code{network/client/Client.java}               & Singleton hub del client \\
\code{network/client/ServerConnection.java}     & Interfaccia astrazione trasporto client \\
\code{network/client/SocketServerConnection.java} & Implementazione TCP client \\
\code{network/client/RMIServerConnection.java}  & Client RMI + target callback \\
\code{network/client/ServerHandler.java}        & Thread lettura background TCP \\
\code{network/client/ClientModel.java}          & Stato locale lato client \\
\code{network/client/VirtualServer.java}        & Facciata azioni per UI \\
\code{network/messages/ClientMessage.java}      & Interfaccia marcatore radice \\
\code{network/messages/Request.java}            & Interfaccia sealed Request \\
\code{network/messages/Update.java}             & Interfaccia sealed Update \\
\code{network/rmi/RMIGameServer.java}           & Interfaccia remota d'ingresso RMI \\
\code{network/rmi/RMIClientInterface.java}      & Interfaccia callback client RMI \\
\code{network/rmi/RMIClientHandlerInterface.java} & Interfaccia handler server RMI \\
\bottomrule
\end{tabularx}
\end{small}

\end{document}
